\chapter{Conclusions}
\label{chapter:conclusions}

This thesis has presented the design, implementation, and evaluation of a prototype data anonymization service for the MECOIS framework, addressing the critical tension \uline{In Einleitung ok, hier too much} between software engineering productivity analytics and individual privacy rights under the \ac{GDPR} and \ac{BDSG}. By integrating this service as a decoupled layer within the MECOIS architecture, the project demonstrates a viable path for organizations to derive actionable insights from developer metadata while maintaining a privacy-first approach.

\section{Summary of Results}

The developed prototype successfully fulfills the core functional and nonfunctional requirements established for a \ac{GDPR} and \ac{BDSG}-compliant benchmarking system. Key achievements of this work include:

\paragraph{Architectural Independence} The service operates as a stateless, containerized system using Docker Compose \uline{ich würde das abstrakter beschreiben als mit dem Produktnamen}, allowing for independent deployment and horizontal scalability without interfering with the primary MECOIS analytics engine. \uline{Außerdem Wiederverwendbarkeit oder Skalierbarkeit auf andere Projekte}

\paragraph{Privacy Mechanisms} The service establishes a robust baseline for privacy by pseudonymizing direct identifiers into deterministic \acp{UUID} and achieving k-anonymity through the aggregation of contributors into organizational equivalence classes.

\paragraph{Seamless Integration} By acting as a proxy for the SortingHat identity management system and providing dedicated endpoints for identity transformation, the service integrates into the existing Bronze-to-Silver and Silver-to-Silver pipelines with minimal disruption.

\paragraph{Data Integrity} Through the \texttt{utils.converter} module, the service ensures that PySpark DataFrames and their schemas are preserved during HTTP transmission, maintaining the analytical utility of high-dimensional metadata. 
\uline{Ich glaube die Vorteile von Datenintegrität kann man auch bodenständiger beschreiben}


\section{Discussion and Legal Compliance}

From a legal perspective, the prototype bridges the gap between scientific research and final product requirements. \uline{Der Satz funktioniert nicht richtig, scientific research bedeutet nicht dass man die rechtlichen Gesichtspunkte ignoriert} While the current use of pseudonymization is sufficient for the research phase of MECOIS, the service provides the necessary \uline{(foundation to build an)} infrastructure to transition toward factual anonymization —  \uline{AI dash} where re-identification requires disproportionate effort — by generalized grouping and the potential permanent deletion of re-identification keys. The evaluation confirms that the system effectively mitigates the risk of "singling out" individuals through its implementation of equivalence classes based on organizational enrollment.

\section{Future Work}

Despite the successful implementation of the prototype, several areas for future development remain to enhance the service's security, privacy robustness, and production readiness.

\subsubsection{Assuring k-anonymity}
While the current prototype achieves a baseline for k-anonymity by grouping contributors into organizational equivalence classes, it lacks an automated mechanism to enforce the k threshold. Future iterations must implement a check to ensure that every organization in a dataset contains at least k enrolled individuals; \uline{Warum kein Punkt?} if this threshold is not met, the system should automatically aggregate smaller organizations into broader clusters to maintain the necessary anonymity level. Furthermore, the system must account for temporal risks such as employee vacations; if a query targets a narrow timeframe where several contributors are inactive, the number of actual records may fall below the k limit even if the total enrollment is sufficient. Implementing dynamic checks for active record counts within queried intervals is essential to prevent re-identification through "singling out".

\subsubsection{Permission System and Logging}
The current service acts as a proxy for the SortingHat identity management system, but it lacks a formal access control layer. Currently, any component within the network can query the system to resolve identities. To align with \ac{GDPR} accountability standards, a robust permission system and authentication layer must be implemented to restrict who can resolve \acp{UUID} or access sensitive profile metadata. Additionally, every query to the identity management system should be logged to provide a transparent audit trail of how and when personal data is processed.

\subsubsection{Production-Grade Server}
The prototype utilizes Python's built-in \texttt{http.server} module to minimize dependencies during development. However, this module is not recommended for production environments as it only implements basic security checks and lacks the performance capabilities for high-volume data pipelines. Future versions should replace this with a production-grade framework such as Flask combined with Gunicorn. This upgrade would provide enhanced security, support for multiple threads, and asynchronous execution, ensuring the service can scale horizontally to handle heavy processing loads behind a load balancer. \autocite{python_http_server_2026}

\subsubsection{Anonymization Analysis}
A critical next step is a formal, comprehensive analysis of the anonymization effectiveness. While the prototype implements key privacy strategies, \uline{ich würde hier auch noch erwähnen dass die Arbeit das auch schon teilsweise evaluiert, aber halt nicht komplett} a complete report is needed to document the "reasonability" of the approach — \uline{AI dash} evaluating whether re-identification requires a disproportionate effort in terms of time and cost. This analysis should be formalized as an "attacker model" that quantifies residual risks such as linkage or inference attacks based on the specific developer metadata processed by MECOIS.

\subsubsection{Cryptographic Security}
The service currently relies on SortingHat's native \ac{UUID} generation, which uses a SHA-1 hashing mechanism. As noted in the evaluation, SHA-1 is susceptible to hash collisions and lacks resistance to brute-force or rainbow table attacks. To improve cryptographic security, future iterations must implement a more secure, salted hashing scheme. Appending a "salt" — a random sequence of characters — before hashing will significantly increase the entropy of the identifiers and prevent attackers from using precomputed tables to reverse the transformation. \uline{eventuell Quelle? Und wenn du das schreibst kannst du auch noch SHA256 schreiben}

\subsubsection{Unstructured Data Processing}
While the service successfully handles structured identifiers, indirect identifiers within unstructured data remains a significant challenge. Specifically, Git commit messages often contain timestamps, specific URLs, or linguistic patterns that could lead to re-identification. Future work should integrate advanced \ac{NLP} for content-level masking or author obfuscation techniques, such as \ac{A4NT}, to mask unique writing styles while preserving the analytical utility of the messages.

\subsubsection{Right to Erasure and Deletion Cascade}
To ensure full compliance with Art. 17 \ac{GDPR}, commonly referred to as the "Right to be Forgotten," the service must implement a mechanism capable of executing a deletion cascade upon a contributor's request. This process requires the system to irreversibly purge or "hard-anonymize" all historical mappings, pseudonym keys, and logs that allow for the re-identification of that specific individual. While the legal equivalence of anonymization and deletion is a subject of academic debate, a robust transformation that renders re-identification impossible—or possible only through a disproportionate effort in time and cost — can satisfy the requirement to make personal data unrecognizable. Within the SortingHat or similar identity management systems, this necessitates the permanent destruction of deterministic \ac{UUID} mapping tables to break the link between the pseudonym and the original direct identifiers. Furthermore, achieving absolute anonymity requires that no party, including the original controller, retains the means to resolve these identifiers. Transitioning the research prototype into a legally compliant product thus \uline{(therefore)} depends on the implementation of these irreversible erasure protocols to uphold the accountability principle and protect the data subject's rights.
\autocite[71--72,126--130]{Krebs2022}
\autocite[13,18--20,73--75]{gmdsbvd2024praxishilfe}
\autocite[13--17]{schwartmann2022praxisleitfaden}
\uline{drei Quellen am Ende vom Textblock würde ich nur machen wenns absolut nicht anders geht}

\subsubsection{Employee Representation}

Various supervisory bodies exist to safeguard and defend employee rights. \uline{Weiterer Satz der sagt dass das mit denen koordiniert werden muss}



\paragraph{Data Protection Officer}

Under certain statutory conditions, organizations are obligated to designate a \ac{DPO}. The \ac{DPO}'s responsibilities encompass providing advisory support, monitoring compliance, and liaising with the supervisory authority. Furthermore, the \ac{DPO} must inform and advise controllers of their legal obligations under data protection law and audit compliance with statutory mandates, such as the \ac{GDPR}. Consequently, the \ac{DPO} must be integrated at an early stage during the planning and design of future digital forms of work \autocite[392--394,400--402]{daeubler2017glaeserne}.

\paragraph{Works Council}

Future implementation of MECOIS in a corporate environment requires the mandatory and active involvement of the works council (\textit{Betriebsrat}), as productivity benchmarking directly affects employee interests. Under German law, specifically \S~87 Abs.~1 Nr.~6 BetrVG, the works council possesses an enforceable right of co-determination regarding the introduction and use of technical facilities designed to monitor the behavior or performance of employees. Since MECOIS processes granular activity-based artifacts such as timestamps and commit details, it falls under this provision regardless of whether the monitoring is the primary intent.
\autocite[429]{daeubler2017glaeserne}
\autocite[2--6,79--81]{niedenhoff2024betriebsraete}


The employer is legally obligated to inform the works council timely and comprehensively  under \S~80 Abs.2 BetrVG and must provide all necessary documentation to allow the council to evaluate the impact on the workforce. Furthermore, if the system utilizes AI for its statistical calculations or composite index generation, the works council has an explicit claim to involve external experts to assess the technical implications under \S80 Abs.~3 BetrVG. To facilitate a legally compliant deployment, a shop agreement (\textit{Betriebsvereinbarung}) must be concluded, which defines the specific boundaries of data usage. The technical measures implemented in this prototype — specifically pseudonymization via deterministic \acp{UUID} and the achievement of k-anonymity through organizational grouping—serve as a critical baseline for these negotiations to ensure that the pursuit of team productivity does not lead to the creation of "glass employees". \uline{(does not come at the expense of workers relations)}
\autocite[2--4,61--62,72--73,96--97,99--101,134,348--351]{niedenhoff2024betriebsraete}
\uline{Der Employee Representation part ist sehr lang und passt thematisch nicht richtig zu den anderen. Ich würde ihn auf jeden Fall nicht als letzten schreiben und ggf. kürzen oder als Abgrenzung oder Definitionen in ein anderes Kapitel schieben}


\vspace{2cm}

In conclusion, the MECOIS anonymization service provides a scalable and legally grounded foundation for modern software analytics, ensuring that the pursuit of team productivity does not come at the expense of developer privacy.
\uline{Der Absatz und danach dieser kurze Textstummel danach ist kein guter Abschluss, das ist zu plötzlich. Wir hatten als letzten Abschnitt unserer Arbeit nach der future Work einen Ausblick entweder aufs Projekt oder die erneute Bezugnahme
auf die Einleitung, in deinem Fall also die Verbeitung von Big Data in allen Branchen und dass du hier deinen Beitrag leistest das besser und privater zu gestalten. Das sollten mindestens fünf Sätze sein mMn.}

